\begin{dang}{BÀI TOÁN QUÃNG ĐƯỜNG}
\setcounter{paragraph}{0}
\paragraph{Quãng đường vật đi trong thời gian $\Delta t$}
\begin{itemize}
\item Phân tích: $\Delta t=nT+\Delta t'\ \text{với n}\ \in \mathbb{N}^*;\ \Delta t'<T$.
\item Tính quãng đường $S_1$ vật đi trong $nT$ với $S_1 = n.4A$.
\item Tính quãng đường $S_2$ vật đi trong $\Delta t'$ bằng một trong hai cách sau:
\begin{itemize}
\item \textbf{Cách 1:} Dùng trục thời gian để xác định quãng đường dịch chuyển từ trạng thái 1 đến trạng thái 2.
\item \textbf{Cách 2:} Dùng vòng tròn lượng giác để xác định quãng đường dịch chuyển từ trạng thái 1 đến trạng thái 2.  
\end{itemize}
\item Quãng đường đi trong thời gian $\Delta t$ là $S=S_1+S_2$
\end{itemize}
\begin{note}
\begin{itemize}
\item Quãng đường vật đi được trong một chu kì $T$ là 4A.
\item Quãng đường vật đi được trong nửa chu kì $\dfrac{T}{2}$ là 2A.
\end{itemize}
\end{note}
\paragraph{Thời gian vật đi được quãng đường cho trước}
\begin{itemize}
\item Đây là bài toán ngược với cách làm hoàn toàn tương tự.
\item Xác định vị trí xuất phát trên VTLG hoặc trục phân bố thời gian.  
\begin{itemize}
\item Nếu $S < 4A$ thì $\Delta t < T$. 
\item  Nếu $S > 4A$ thì $\Delta t > T$ :
$S=\underbrace{n.4A}_{nT}+\underbrace{S_{\text{thêm}}}_{\Delta t'}\Rightarrow \Delta t=nT+\Delta t'$.  
\end{itemize}
\item Phối hợp vòng tròn lượng giác với trục thời gian để xác định khoảng thời gian.
\end{itemize}
\end{dang}